\chapter{Conclusion générale et perspectives}
\label{chap:conclusion}

\section{Conclusion générale}
\label{sec:conclusion:generale}

Le projet \projectname{} nous a permis de concevoir et réaliser une plateforme complète de génération de CV assistée par intelligence artificielle. Ce projet d'intégration a mobilisé l'ensemble des compétences acquises durant notre formation, du développement web à l'architecture microservices, en passant par la conteneurisation, l'intégration continue, la sécurité applicative et l'observabilité.

L'approche DevOps adoptée, avec une organisation en microservices et un pipeline CI/CD complet intégrant qualité et sécurité, a permis de livrer une application robuste et industrialisée. La pipeline d'agents IA basée sur LangChain et LangGraph démontre la capacité à intégrer des technologies avancées d'intelligence artificielle dans une application concrète.

Sur le plan méthodologique, l'adoption de SCRUM avec des réunions quotidiennes et une communication continue via Google Meets et WhatsApp a favorisé une collaboration efficace au sein de l'équipe de six développeurs fullstack, chacun apportant son expertise dans son domaine de responsabilité secondaire.

\section{Perspectives}
\label{sec:conclusion:perspectives}

Plusieurs perspectives d'amélioration et d'évolution sont envisageables pour \projectname{} :

\begin{itemize}
    \item \textbf{Déploiement cloud :} Migrer l'infrastructure vers un fournisseur cloud (AWS, Azure, GCP) avec orchestration Kubernetes pour une meilleure scalabilité.
    \item \textbf{Modèles de langage additionnels :} Intégrer d'autres fournisseurs de LLM (OpenAI, Anthropic, modèles open source en local) pour diversifier les capacités de génération.
    \item \textbf{Internationalisation :} Ajouter le support multilingue pour l'interface et la génération de CV dans différentes langues.
    \item \textbf{Analyse avancée :} Développer des fonctionnalités d'analyse prédictive pour suggérer les offres d'emploi les plus pertinentes en fonction du profil de l'utilisateur.
    \item \textbf{Mode hors ligne :} Proposer une version allégée de l'application fonctionnant sans connexion Internet pour la consultation des CV et candidatures.
    \item \textbf{Place de marché :} Créer une plateforme de mise en relation entre candidats et recruteurs utilisant les données de profiling générées par l'application.
\end{itemize}
